%auto-ignore
%this ensures the arxiv doesn't try to start TeXing here.
% !TEX root = exceptional.tex

\usepackage{fp}
\usepackage{tikz}
\usetikzlibrary{matrix}
\usetikzlibrary{arrows,backgrounds,patterns,scopes,external,hobby,
    decorations.pathreplacing,
    decorations.pathmorphing
}
\usepackage{tikzit}


\newlength{\fuzzwidth}
\setlength{\fuzzwidth}{2.5pt}
\newlength{\arrowlength}
\setlength{\arrowlength}{8pt}
\newlength{\arrowwidth}
\setlength{\arrowwidth}{.75pt}
\newlength{\pointrad}
\setlength{\pointrad}{1.5pt}
\newlength{\linewid}
\setlength{\linewid}{1.5pt}
\newlength{\circlerad}
\setlength{\circlerad}{16pt}
\newlength{\smcirclerad}
\setlength{\smcirclerad}{8pt}
\newcommand{\fillcolor}{black!60}
\newcommand{\fuzzcolor}{black!25}
\newcommand{\arrowcolor}{black!25}
\newcommand{\covercolor}{black!0}
\newcommand{\graycolor}{black!55}
\newcommand{\graylightcolor}{black!40}


\newcommand{\coverwidthfuzz}{6pt}
\newcommand{\coverwidth}{3.5pt}
\newcommand{\coverwidththin}{3.25pt}
\newcommand{\coverwidththick}{3.75pt}

\newlength{\linewidthin}
\setlength{\linewidthin}{1.25pt}
\newlength{\linewidthick}
\setlength{\linewidthick}{1.75pt}

\tikzset{cdlabel/.style={execute at begin node=$\scriptstyle,execute at end node=$}}

\tikzset{use Hobby shortcut}
\tikzset{
	coverline/.style={
	preaction={draw,line width=\coverwidth,\covercolor}}, 
	coverlinethin/.style={
	preaction={draw,line width=\coverwidththin,\covercolor}}, 
	coverlinethick/.style={
	preaction={draw,line width=\coverwidththick,\covercolor}}, 
	coverlineleft/.style={
	preaction={draw,line width=\coverwidthfuzz,\covercolor,decorate,decoration={curveto,amplitude=0,raise=.35*\fuzzwidth}}}, %%% Notice this number was tweaked, 	
coverlinelefttail/.style={
	preaction={draw,line width=\coverwidthfuzz,\covercolor,decorate,decoration={curveto,amplitude=0,raise=.35*\fuzzwidth,pre=moveto,pre length=2pt}}}, %%% Notice this number was tweaked, probably will not look good if parameters change.  Again here, pre removes the raise option.
        fuzzlefttail/.style={
        preaction={draw,line width=\fuzzwidth,\fuzzcolor,decorate,decoration={curveto,pre=moveto,pre length=2pt,amplitude=0,raise=.5*\fuzzwidth}}}, %%% This does not work --- somehow pre breaks the raise function.
        linestylethin/.style={line width=\linewidthin},
        linestylethick/.style={line width=\linewidthick},
        linestylegray/.style={line width=\linewid,\graycolor},
        linestylegraylight/.style={line width=\linewid,\graylightcolor}
}
\tikzset{
        fuzzright/.style={
        preaction={draw,line width=\fuzzwidth,\fuzzcolor,decorate,decoration={curveto,amplitude=0,raise=-.5*\fuzzwidth}}},
        fuzzleft/.style={
        preaction={draw,line width=\fuzzwidth,\fuzzcolor,decorate,decoration={curveto,amplitude=0,raise=.5*\fuzzwidth}}},
        fuzzrightpre/.style={ %%% Doesn't work
        preaction={draw,line width=2pt,\fuzzcolor,decorate,decoration={curveto,amplitude=0,raise=-1pt,pre=moveto,pre length=12pt}}},
        fuzzleftpre/.style={ %%% Doesn't work
        preaction={draw,line width=2pt,\fuzzcolor,decorate,decoration={curveto,post=moveto,post length=32pt,amplitude=0,raise=1pt}}},        
        outstyle/.style={\arrowcolor, line width=\arrowwidth},
        linestyle/.style={line width=\linewid}
}

\newcommand{\lilarrow}{
\draw[->] (0,0) -- (1,0);
}

\newcommand{\cb}{\raisebox{.6ex-.5\height}}


%%%% These draw triple or quadruple set of arrows of length 0.5 cm
\DeclareMathOperator{\rightdoublearrows} {{\; \begin{tikzpicture} \foreach \y in {0.05, 0.15} { \draw [-stealth] (0, \y) -- +(0.5, 0);} \; \end{tikzpicture}}}
\DeclareMathOperator{\leftdoublearrows} {{\; \begin{tikzpicture} \foreach \y in {0.05, 0.15} { \draw [stealth-] (0, \y) -- +(0.5, 0);} \; \end{tikzpicture}}}
\DeclareMathOperator{\righttriplearrows} {{\; \begin{tikzpicture} \foreach \y in {0, 0.1, 0.2} { \draw [-stealth] (0, \y) -- +(0.5, 0);} \; \end{tikzpicture}}}
\DeclareMathOperator{\lefttriplearrows} {{\; \begin{tikzpicture} \foreach \y in {0, 0.1, 0.2} { \draw [stealth-] (0, \y) -- +(0.5, 0);} \; \end{tikzpicture}}}
\DeclareMathOperator{\rightquadarrows} {{\; \begin{tikzpicture} \foreach \y in {-0.05, 0.05, 0.15, 0.25} { \draw [-stealth] (0, \y) -- +(0.5, 0);} \; \end{tikzpicture}}}
\DeclareMathOperator{\leftquadarrows} {{\; \begin{tikzpicture} \foreach \y in {-0.05, 0.05, 0.15, 0.25} { \draw [stealth-] (0, \y) -- +(0.5, 0);} \; \end{tikzpicture}}}

\newcommand{\newfig}[3]{
  \newsavebox{#2}
  \savebox{#2}{#3}
  \newcommand{#1}{\usebox{#2}}
}

\newcommand{\smallfig}[1]{\mathcenter{\scalebox{0.5}{#1}}}

\newfig{\twistedsquarehor}{\twistedsquarehorbox}{
\begin{tikzpicture}[baseline=-.5ex,scale=.8]
\draw (45:.8cm) -- (-135:.8cm);
\draw[line width=1mm,white,double=black] (-45:.8cm) -- (135:.8cm);
\draw (45:.5cm) .. controls +(-75:.3cm) and +(75:.3cm) .. (-45:.5cm);
\draw (135:.5cm) .. controls +(-105:.3cm) and +(105:.3cm) .. (-135:.5cm);
\end{tikzpicture}}

\newfig{\twistedsquarever}{\twistedsquareverbox}{
\begin{tikzpicture}[baseline=-.5ex,scale=.8]
\draw (45:.8cm) -- (-135:.8cm);
\draw[line width=1mm,white,double=black] (-45:.8cm) -- (135:.8cm);
\draw (45:.5cm) .. controls +(165:.3cm) and +(15:.3cm) .. (135:.5cm);
\draw (-45:.5cm) .. controls +(-165:.3cm) and +(-15:.3cm) .. (-135:.5cm);
\end{tikzpicture}}



\newcommand{\ngon}[2][0]{
\begin{tikzpicture}[baseline=-0.5ex,scale=0.8]
\foreach \x in {1, ..., #2}
	\draw (360*\x/#2+#1:.8cm)--(360*\x/#2+#1:.5cm);
\foreach \x in {1, ..., #2}
	\draw (360*\x/#2+#1:.5cm) .. controls +(360*\x/#2+120+#1:.3cm) and +(360*\x/#2+360/#2-120+#1:.3cm) .. (360*\x/#2+360/#2+#1:.5cm);
\end{tikzpicture}
}

\newcommand{\nvertex}[2][0]{
\begin{tikzpicture}[baseline=-0.5ex,scale=0.8]
\foreach \x in {1, ..., #2}
	\draw (360*\x/#2+#1:.8cm)--(0,0);
\end{tikzpicture}
}

\newfig{\twogon}{\twogonbox}{\ngon{2}}
\newfig{\threegon}{\threegonbox}{\ngon[-90]{3}}
\newfig{\threevertex}{\threevertexbox}{\nvertex[-90]{3}}
\newfig{\fourgon}{\fourgonbox}{\ngon[45]{4}}

\newfig{\upsidedownthreevertex}{\upsidedownthreevertexbox}{%
\begin{tikzpicture}[baseline=-0.5ex,scale=0.8]
	\draw (330:.8cm)--(0,0);
	\draw (90:.8cm)--(0,0);
	\draw (210:.8cm)--(0,0);
\end{tikzpicture}}

\newfig{\unknot}{\unknotbox}{%
\begin{tikzpicture}[baseline=-0.5ex,scale=0.8]
  \draw (0,0) circle (.6cm);
\end{tikzpicture}}

\newfig{\drawI}{\drawIbox}{\begin{tikzpicture}[baseline=-0.5ex,scale=0.8]
 	\draw (0,.2) -- (45:.8cm);
 	\draw (0,.2) -- (135:.8cm);
	\draw (0,.2) -- (0,-.2);
 	\draw (0,-.2) -- (-45:.8cm);
 	\draw (0,-.2) -- (-135:.8cm);
\end{tikzpicture}}

\newfig{\drawH}{\drawHbox}{\begin{tikzpicture}[baseline=-0.5ex,rotate=90,scale=0.8]
 	\draw (0,.2) -- (45:.8cm);
 	\draw (0,.2) -- (135:.8cm);
	\draw (0,.2) -- (0,-.2);
 	\draw (0,-.2) -- (-45:.8cm);
 	\draw (0,-.2) -- (-135:.8cm);
\end{tikzpicture}}

\newfig{\drawdottedH}{\drawdottedHbox}{\begin{tikzpicture}[baseline=-0.5ex,scale=0.8]
        \draw (45:.8cm) -- (-45:.8cm);
        \draw (135:.8cm) -- (-135:.8cm);
        \draw[dotted] (-0.57cm,0) -- (0.57cm, 0);
\end{tikzpicture}}

\newfig{\onestrandid}{\onestrandidbox}{\begin{tikzpicture}[baseline=-0.5ex,scale=0.8]
	\draw (-.8cm,0)--(.8cm,0);
\end{tikzpicture}}

\newfig{\vertonestrandid}{\vertonestrandidbox}{\begin{tikzpicture}[rotate=90,baseline=-0.5ex,scale=0.8]
	\draw (-.8cm,0)--(.8cm,0);
\end{tikzpicture}}


\newfig{\drawcup}{\drawcupbox}{\begin{tikzpicture}[baseline=-0.5ex,scale=0.8]
	\draw (45:.8cm) to [curve through=(90:0cm)] (135:.8cm);
\end{tikzpicture}}

\newfig{\drawcap}{\drawcapbox}{\begin{tikzpicture}[baseline=-0.5ex,scale=0.8]
	\draw (-45:.8cm) to [curve through=(-90:0cm)] (-135:.8cm);
\end{tikzpicture}}

\newfig{\drawbubblecap}{\drawbubblecapbox}{\begin{tikzpicture}[baseline=-0.5ex,scale=0.8]
	\draw (180:.8cm) to [curve through = (135: .8cm)] (120:.8cm);
	\draw (0:.8cm) to [curve through = (45: .8cm)] (60:.8cm);
	\draw (120:.8cm) to [curve through = (90:1cm)] (60:.8cm);
	\draw (120:.8cm) to [curve through = (90:.4cm)] (60:.8cm);
\end{tikzpicture}}



\newfig{\cupcap}{\cupcapbox}{\begin{tikzpicture}[baseline=-0.5ex,scale=0.8]
	\draw (45:.8cm) to [curve through=(90:.3cm)] (135:.8cm);
	\draw (-45:.8cm) to [curve through=(-90:.3cm)] (-135:.8cm);
\end{tikzpicture}}

\newfig{\twostrandid}{\twostrandidbox}{\begin{tikzpicture}[baseline=-0.5ex,rotate=90,scale=0.8]
	\draw (45:.8cm) to [curve through=(90:.3cm)] (135:.8cm);
	\draw (-45:.8cm) to [curve through=(-90:.3cm)] (-135:.8cm);
\end{tikzpicture}}

\newfig{\symcross}{\symcrossbox}{\begin{tikzpicture}[baseline=-0.5ex,scale=0.8]
	\draw (45:.8cm) -- (-135:.8cm);
	\draw (-45:.8cm) -- (135:.8cm);
\end{tikzpicture}}

\newfig{\braidcross}{\braidcrossbox}{\begin{tikzpicture}[baseline=-0.5ex,scale=0.8]
	\draw (45:.8cm) -- (-135:.8cm);
	\draw[line width=1mm,white,double=black] (-45:.8cm) -- (135:.8cm);
\end{tikzpicture}}

\newfig{\invbraidcross}{\invbraidcrossbox}{\begin{tikzpicture}[baseline=-0.5ex,scale=0.8]
	\draw (-45:.8cm) -- (135:.8cm);
	\draw[line width=1mm,white,double=black] (45:.8cm) -- (-135:.8cm);
\end{tikzpicture}}


\newfig{\drawcrossX}{\drawcrossXbox}{\begin{tikzpicture}[baseline=-0.5ex,scale=0.8,rotate=90]
        \draw ([out angle=-150].2,-.3) ..([in angle=-70]135:.8cm);
 	\draw[line width=2mm,white]
              ([out angle=150].2,.3) .. ([in angle=70]-135:.8cm);
 	\draw ([out angle=150].2,.3) .. ([in angle=70]-135:.8cm);
 	\draw ([out angle=30].2,.3) .. (45:.8cm);
	\draw (.2,.3) -- (.2,-.3);
 	\draw ([out angle=-30].2,-.3) .. (-45:.8cm);
\end{tikzpicture}}

\newfig{\drawinvcrossX}{\drawinvcrossXbox}{\begin{tikzpicture}[baseline=-0.5ex,scale=0.8,rotate=90]
 	\draw ([out angle=150].2,.3) .. ([in angle=70]-135:.8cm);
        \draw[line width=2mm,white] ([out angle=-150].2,-.3) ..([in angle=-70]135:.8cm);
        \draw([out angle=-150].2,-.3) ..([in angle=-70]135:.8cm);
	\draw ([out angle=30].2,.3) .. (45:.8cm);
	\draw (.2,.3) -- (.2,-.3);
 	\draw ([out angle=-30].2,-.3) .. (-45:.8cm);
\end{tikzpicture}}


\newfig{\drawsymcrossX}{\drawsymcrossXbox}{\begin{tikzpicture}[baseline=-0.5ex,scale=0.8,rotate=90]
 	\draw ([out angle=30].2,.3) .. (45:.8cm);
	\draw (.2,.3) -- (.2,-.3);
 	\draw ([out angle=-30].2,-.3) .. (-45:.8cm);
        \draw ([out angle=-150].2,-.3) ..([in angle=-70]135:.8cm);
 	\draw%[line width=1mm,white,double=black]
              ([out angle=150].2,.3) .. ([in angle=70]-135:.8cm);
\end{tikzpicture}}

\newfig{\drawsyminvcrossX}{\drawsyminvcrossXbox}{\begin{tikzpicture}[baseline=-0.5ex,scale=0.8,rotate=-90]
 	\draw ([out angle=30].2,.3) .. (45:.8cm);
	\draw (.2,.3) -- (.2,-.3);
 	\draw ([out angle=-30].2,-.3) .. (-45:.8cm);
        \draw ([out angle=-150].2,-.3) ..([in angle=-70]135:.8cm);
 	\draw%[line width=1mm,white,double=black]
              ([out angle=150].2,.3) .. ([in angle=70]-135:.8cm);
\end{tikzpicture}}


\newfig{\twist}{\twistbox}{%
  \begin{tikzpicture}[baseline=-0.5ex,scale=0.8]
    \draw[line width=1mm,white,double=black]
       ([out angle=-15]135:.8cm)..([blank=soft]-60:.4cm)..(-120:.4cm)..([in angle=-165]45:.8cm);
    \draw[line width=1mm,white,double=black,use previous Hobby path={invert soft blanks,disjoint}];
  \end{tikzpicture}}


\newfig{\symtwist}{\symtwistbox}{%
  \begin{tikzpicture}[baseline=-0.5ex,scale=0.8]
    \draw
       ([out angle=-15]135:.8cm)..([blank=soft]-60:.4cm)..(-120:.4cm)..([in angle=-165]45:.8cm);
   \draw[use previous Hobby path={invert soft blanks,disjoint}];
  \end{tikzpicture}}


\newfig{\twistvertex}{\twistvertexbox}{%
\begin{tikzpicture}[baseline=-0.5ex,scale=0.8]
  \draw (0,-0.4)--(0,-0.8);
  \draw ([out angle=-170]30:0.8cm)..([in angle=160](0,-0.4);
  \draw[line width=1mm,white,double=black]
        ([out angle=-10]150:0.8cm)..([in angle=20](0,-0.4);
\end{tikzpicture}}

\newfig{\symtwistvertex}{\symtwistvertexbox}{%
\begin{tikzpicture}[baseline=-0.5ex,scale=0.8]
  \draw (0,-0.4)--(0,-0.8);
  \draw ([out angle=-170]30:0.8cm)..([in angle=160](0,-0.4);
  \draw%[line width=1mm,white,double=black]
        ([out angle=-10]150:0.8cm)..([in angle=20](0,-0.4);
\end{tikzpicture}}


\newfig{\loopvertex}{\loopvertexbox}{%
\begin{tikzpicture}[baseline=-0.5ex,scale=0.8]
  \draw (0,-0.5)--(0,-0.8);
  \draw ([out angle=150]0,-0.5)..(-0.02,0.6)..(0.02,0.6)..([in angle=30]0,-0.5);
\end{tikzpicture}}

\newfig{\idtangle}{\idtanglebox}{%
\begin{tikzpicture}[baseline=-0.5ex,scale=0.8]
\draw (0,0) node {$D$};
\draw[densely dashed] (0,0) circle (.5cm);
\draw (45:.5cm) -- (45:0.8cm);
\draw (135:.5cm) -- (135:0.8cm);
\draw (-45:.5cm) -- (-45:0.8cm);
\draw (-135:.5cm) -- (-135:0.8cm);
\end{tikzpicture}}
    
\newfig{\Rcycle}{\Rcyclebox}{%
\begin{tikzpicture}[baseline=-0.5ex,scale=0.8]
\draw (0,0) node {$D$};
\draw[densely dashed] (0,0) circle (.4cm);
\draw (45:.4cm) -- (45:1.1cm);
\draw ([out angle=-70]-45:.4cm) .. (-90:.7cm) .. ([in angle=0]-135:1.1cm);
\draw ([out angle=-160]-135:.4cm) .. (180:.7cm) .. ([in angle=-90]135:1.1cm);
\draw[line width=0.8mm,white,double=black]
      ([out angle=90]135:.4cm) .. (45:.7cm) .. ([in angle=90]-45:1.1cm);
\end{tikzpicture}}

\newfig{\ladder}{\ladderbox}{%
\begin{tikzpicture}[baseline=-0.5ex,scale=0.8]
\draw (0,0) node {$D$};
\draw[densely dashed] (0,0) circle (.5cm);
\draw (135:.5cm) -- (135:0.9cm);
\draw (-135:.5cm) -- (-135:0.9cm);
\draw ([out angle=45]45:.5cm) .. ([in angle=150]0.8,0.45);
\draw ([out angle=-45]-45:.5cm) .. ([in angle=-150]0.8,-0.45);
\draw (0.8,0.45)--(0.8,-0.45);
\draw ([out angle=30]0.8,0.45) .. ([in angle=-135]1.2,0.636);
\draw ([out angle=-30]0.8,-0.45) .. ([in angle=135]1.2,-0.636);
\end{tikzpicture}}

% \newfig{\pentagon}{\pentagonbox}{
% \begin{tikzpicture}[scale=.12,baseline=-2]
% \draw (36:1) -- (108:1) -- (180:1) -- (252:1) -- (-36:1) -- (36:1);
% \end{tikzpicture}
% }

% \newfig{\psq}{\psqbox}{
% \begin{tikzpicture}[scale=.15,baseline=-2]
% \draw (36:1) -- (108:1) -- (180:1) -- (252:1) -- (-36:1) -- (36:1) -- +(1.2,0) -- ($(-36:1)+(1.2,0)$) -- (-36:1);
% \end{tikzpicture}
% }

\newfig{\sqsq}{\sqsqbox}{%
\begin{tikzpicture}[scale=.15]
\draw (0,0) -- (2,0) -- (2,1) -- (0,1) -- cycle (1,0) -- (1,1);
\end{tikzpicture}}

\newfig{\overviolin}{\overviolinbox}{%
\begin{tikzpicture}[baseline=0]
	\coordinate (E1) at (45:.8cm);
	\coordinate (E2) at (225:.8cm);
	\coordinate (E3) at (270:.8cm);
	\coordinate (P1) at (0,0);
	\coordinate (P2) at (120:.6cm);
	\coordinate (P3) at (.4cm,.2cm);
	\coordinate (C1) at (100:.9cm);
	\coordinate (C2) at (135:.5cm);
	\coordinate (C3) at (.8,.1);
	\draw (E1) -- (P1) -- (E3);
	\draw (P1) .. controls  (C1) .. (P2);
	\draw[white, line width=4pt] (P2) .. controls  (C2) .. (P3);	
	\draw[white, line width=4pt] (P3) .. controls (C3) .. (E2);
	\draw (P2) .. controls  (C2) .. (P3);	
	\draw (P3) .. controls (C3) .. (E2);
%	\path (P1) ++(330:.15cm) node {$\bullet$};
\end{tikzpicture}}

\newfig{\rotatedtrivalent}{\rotatedtrivalentbox}{%
\begin{tikzpicture}[baseline=0]
	\coordinate (E1) at (45:.8cm);
	\coordinate (E2) at (225:.8cm);
	\coordinate (E3) at (270:.8cm);
	\coordinate (P0) at (0,0);
	\coordinate (C1) at (-.1,.7);
	\coordinate (C2) at (-.3,-.2);
	\coordinate (C3) at (.2,-.2);
	\draw (P0) .. controls (C1) .. (E2);
	\draw (P0) .. controls (C2) .. (E3);
	\draw (P0) .. controls (C3) .. (E1);
%	\path (P0) ++(270:.2cm) node {$\bullet$};
\end{tikzpicture}}

\newcommand{\undercrossingfork}{%
\begin{tikzpicture}[baseline=0]
	\coordinate (E1) at (90-72:.8cm);
	\coordinate (E2) at (90:.8cm);
	\coordinate (E3) at (90+72:.8cm);
	\coordinate (E4) at (90+2*72:.8cm);
	\coordinate (E5) at (90+3*72:.8cm);
	\coordinate (P1) at (270:.2cm);
	\draw (E4) -- (P1) -- (E5);
	\draw (P1) -- (E2);
	\draw[white, line width = 4pt] (E1) -- (E3);
	\draw (E1) -- (E3);
%	\path (P1) ++(270:.2cm) node {$\bullet$};
\end{tikzpicture}}

\newcommand{\undercrossingforktwoleft}{%
\begin{tikzpicture}[baseline=0,rotate =144]
	\coordinate (E1) at (90-72:.8cm);
	\coordinate (E2) at (90:.8cm);
	\coordinate (E3) at (90+72:.8cm);
	\coordinate (E4) at (90+2*72:.8cm);
	\coordinate (E5) at (90+3*72:.8cm);
	\coordinate (P1) at (270:.2cm);
	\draw (E4) -- (P1) -- (E5);
	\draw (P1) -- (E2);
	\draw[white, line width = 4pt] (E1) -- (E3);
	\draw (E1) -- (E3);
%	\path (P1) ++(270:.2cm) node {$\bullet$};
\end{tikzpicture}}


\newcommand{\overcrossingfork}{%
\begin{tikzpicture}[baseline=0]
	\coordinate (E1) at (90-72:.8cm);
	\coordinate (E2) at (90:.8cm);
	\coordinate (E3) at (90+72:.8cm);
	\coordinate (E4) at (90+2*72:.8cm);
	\coordinate (E5) at (90+3*72:.8cm);
	\coordinate (P1) at (270:.2cm);
	\draw (E4) -- (P1) -- (E5);
	\draw (E1) -- (E3);
	\draw[white, line width = 4pt] (P1) -- (E2);
	\draw (P1) -- (E2);
%	\path (P1) ++(270:.2cm) node {$\bullet$};
\end{tikzpicture}}

\newcommand{\overcrossingforktworight}{%
\begin{tikzpicture}[baseline=0, rotate=216]
	\coordinate (E1) at (90-72:.8cm);
	\coordinate (E2) at (90:.8cm);
	\coordinate (E3) at (90+72:.8cm);
	\coordinate (E4) at (90+2*72:.8cm);
	\coordinate (E5) at (90+3*72:.8cm);
	\coordinate (P1) at (270:.2cm);
	\draw (E4) -- (P1) -- (E5);
	\draw (E1) -- (E3);
	\draw[white, line width = 4pt] (P1) -- (E2);
	\draw (P1) -- (E2);
%	\path (P1) ++(270:.2cm) node {$\bullet$};
\end{tikzpicture}}



\newcommand{\underunderfork}{%
\begin{tikzpicture}[baseline=0]
	\coordinate (E1) at (90-72:.8cm);
	\coordinate (E2) at (90:.8cm);
	\coordinate (E3) at (90+72:.8cm);
	\coordinate (E4) at (90+2*72:.8cm);
	\coordinate (E5) at (90+3*72:.8cm);
	\coordinate (P1) at (90:.2cm);
	\draw (E4) .. controls (180:.2cm) .. (P1) .. controls (0:.2cm) ..  (E5);
	\draw (P1) -- (E2);
	\draw[white, line width = 4pt] (E1) .. controls (270:.4cm) .. (E3);
	\draw (E1) .. controls (270:.4cm) .. (E3);
\end{tikzpicture}}

\newcommand{\underoverfork}{%
\begin{tikzpicture}[baseline=0]
	\coordinate (E1) at (90-72:.8cm);
	\coordinate (E2) at (90:.8cm);
	\coordinate (E3) at (90+72:.8cm);
	\coordinate (E4) at (90+2*72:.8cm);
	\coordinate (E5) at (90+3*72:.8cm);
	\coordinate (P1) at (90:.2cm);
	\coordinate (P2) at (90:.-.4cm);
	\draw (P1) .. controls (0:.2cm) ..  (E5);
	\draw (P1) -- (E2);
	\draw[white, line width = 4pt] (E1) .. controls (270:.4cm) .. (E3);
	\draw (E1) .. controls (270:.4cm) .. (E3);
	\draw[white, line width = 4pt] (E4) .. controls (180:.2cm) .. (P1);
	\draw (E4) .. controls (180:.2cm) .. (P1);
\end{tikzpicture}}

\newcommand{\overunderfork}{%
\begin{tikzpicture}[baseline=0]
	\coordinate (E1) at (90-72:.8cm);
	\coordinate (E2) at (90:.8cm);
	\coordinate (E3) at (90+72:.8cm);
	\coordinate (E4) at (90+2*72:.8cm);
	\coordinate (E5) at (90+3*72:.8cm);
	\coordinate (P1) at (90:.2cm);
	\coordinate (P2) at (90:.-.4cm);
	\draw (E4) .. controls (180:.2cm) .. (P1);
	\draw (P1) -- (E2);
	\draw[white, line width = 4pt] (E1) .. controls (270:.4cm) .. (E3);
	\draw (E1) .. controls (270:.4cm) .. (E3);
	\draw[white, line width = 4pt] (P1) .. controls (0:.2cm) ..  (E5);
	\draw (P1) .. controls (0:.2cm) ..  (E5);	
\end{tikzpicture}}


\newcommand{\overoverfork}{%
\begin{tikzpicture}[baseline=0]
	\coordinate (E1) at (90-72:.8cm);
	\coordinate (E2) at (90:.8cm);
	\coordinate (E3) at (90+72:.8cm);
	\coordinate (E4) at (90+2*72:.8cm);
	\coordinate (E5) at (90+3*72:.8cm);
	\coordinate (P1) at (90:.2cm);
	\draw (E1) .. controls (270:.4cm) .. (E3);
	\draw[white, line width = 4pt] (E4) .. controls (180:.2cm) .. (P1) .. controls (0:.2cm) ..  (E5);
	\draw (E4) .. controls (180:.2cm) .. (P1) .. controls (0:.2cm) ..  (E5);
	\draw (P1) -- (E2);
\end{tikzpicture}}

\newcommand{\pentagon}{%
\begin{tikzpicture}[baseline=0, use Hobby shortcut]
	\coordinate (E1) at (90-72:.8cm);
	\coordinate (E2) at (90:.8cm);
	\coordinate (E3) at (90+72:.8cm);
	\coordinate (E4) at (90+2*72:.8cm);
	\coordinate (E5) at (90+3*72:.8cm);
	\coordinate (F1) at (90-72:.4cm);
	\coordinate (F2) at (90:.4cm);
	\coordinate (F3) at (90+72:.4cm);
	\coordinate (F4) at (90+2*72:.4cm);
	\coordinate (F5) at (90+3*72:.4cm);
        \draw (F1)--(F2)--(F3)--(F4)--(F5)--cycle;
        \draw (E1)--(F1);
        \draw (E2)--(F2);
        \draw (E3)--(F3);
        \draw (E4)--(F4);
        \draw (E5)--(F5);
\end{tikzpicture}}

\newcommand{\pentasquare}{%
\begin{tikzpicture}[baseline=0, use Hobby shortcut]
	\coordinate (E1) at (90-72:.8cm);
	\coordinate (E2) at (90:.8cm);
	\coordinate (E3) at (90+72:.8cm);
	\coordinate (E4) at (90+2*72:.8cm);
	\coordinate (E5) at (90+3*72:.8cm);
	\coordinate (F1) at (90-72:.4cm);
	\coordinate (F2) at (90:.4cm);
	\coordinate (F3) at (90+72:.4cm);
	\coordinate (F4) at ($(90+2*72:.3cm)+(0,0.1cm)$);
	\coordinate (F5) at ($(90+3*72:.3cm)+(0,0.1cm)$);
        \coordinate (G4) at ($(F4)-(0,.3cm)$);
        \coordinate (G5) at ($(F5)-(0,.3cm)$);
        \draw (F1)--(F2)--(F3)--(F4)--(F5)--cycle;
        \draw (E1)--(F1);
        \draw (E2)--(F2);
        \draw (E3)--(F3);
        \draw (E4)--(G4)--(F4);
        \draw (E5)--(G5)--(F5);
        \draw (G4)--(G5);
\end{tikzpicture}}

\newcommand{\twisttreebase}{%
	\coordinate (E1) at (90-72:.8cm);
	\coordinate (E2) at (90:.8cm);
	\coordinate (E3) at (90+72:.8cm);
	\coordinate (E4) at (90+2*72:.8cm);
	\coordinate (E5) at (90+3*72:.8cm);
	\coordinate (F1) at (90-72:.4cm);
	\coordinate (F2) at (90:.4cm);
	\coordinate (F3) at (90+72:.4cm);
        \draw (E4)--(F1);
	\draw[white, line width = 4pt] (E5) -- (F3);
        \draw (E5) -- (F3);
        \draw (F1)--(F2)--(F3);
        \draw (E1)--(F1);
        \draw (E2)--(F2);
        \draw (E3)--(F3);
}

\newcommand{\twisttreefourfive}{%
\begin{tikzpicture}[baseline=0, use Hobby shortcut]
\twisttreebase
\end{tikzpicture}}

\newcommand{\twisttreethreefour}{%
\begin{tikzpicture}[baseline=0, use Hobby shortcut,rotate=-72]
\twisttreebase
\end{tikzpicture}}

\newcommand{\twisttreetwothree}{%
\begin{tikzpicture}[baseline=0, use Hobby shortcut,rotate=-2*72]
\twisttreebase
\end{tikzpicture}}

\newcommand{\twisttreeonetwo}{%
\begin{tikzpicture}[baseline=0, use Hobby shortcut,rotate=2*72]
\twisttreebase
\end{tikzpicture}}

\newcommand{\twisttreefiveone}{%
\begin{tikzpicture}[baseline=0, use Hobby shortcut,rotate=72]
\twisttreebase
\end{tikzpicture}}

\newcommand{\twisttreefivefour}{%
\begin{tikzpicture}[baseline=0, use Hobby shortcut,xscale=-1]
\twisttreebase
\end{tikzpicture}}

\newcommand{\twisttreefourthree}{%
\begin{tikzpicture}[baseline=0, use Hobby shortcut,rotate=-72,xscale=-1]
\twisttreebase
\end{tikzpicture}}

\newcommand{\twisttreethreetwo}{%
\begin{tikzpicture}[baseline=0, use Hobby shortcut,rotate=-2*72,xscale=-1]
\twisttreebase
\end{tikzpicture}}

\newcommand{\twisttreetwoone}{%
\begin{tikzpicture}[baseline=0, use Hobby shortcut,rotate=2*72,xscale=-1]
\twisttreebase
\end{tikzpicture}}

\newcommand{\twisttreeonefive}{%
\begin{tikzpicture}[baseline=0, use Hobby shortcut,rotate=72,xscale=-1]
\twisttreebase
\end{tikzpicture}}

\newcommand{\casimirbubble}{%
\begin{tikzpicture}[baseline=-.5ex,scale=0.8]
\draw[->]  (0,-1) -- (0,-.75);
\draw[->]  (0,-.75) -- (0,0);
\draw[->] (0,0) -- (0,.87);
\draw  (0,.87) -- (0,1);
\draw (0,-.5) .. controls (-.5,-.25) and (-.5,.25) .. (0,.5) ;
\end{tikzpicture}}

\newcommand{\orientedonestrandid}{%
\begin{tikzpicture}[baseline=-.5ex,scale=0.8]
\draw[->]  (0,-1) -- (0,0);
\draw  (0,0) -- (0,1);
\end{tikzpicture}}

\newcommand{\ladderaa}{%
\begin{tikzpicture}[baseline=-1.5cm]
		\node[draw, rectangle] (coupon) at (-3.75, -1.5) {$\pi_{V_\mu}$};
		\node  (SW) at (-4, -2.5) {};
		\node  (SE) at (-3.5, -2.5) {};
		\node  (NW) at (-4, -1) {};
		\node  (NE) at (-3.5, -1) {};
		\node  (NofNW) at (-4, 0) {};
		\node  (NofNE) at (-3.5, 0) {};
		\node  (NNofNW) at (-4, .5) {};
		\node  (NNofNE) at (-3.5, .5) {};
		\draw (SW.center) to (coupon.230);
		\draw (SE.center) to (coupon.310);
		\draw (coupon.130) to (NW.center);
		\draw (coupon.50) to (NE.center);
		\draw (NW.center) to (NofNW.center);
		\draw (NE.center) to (NofNE.center);
		\draw [bend left=75, looseness=1] (NW.center) to (NofNW.center);
		\draw (NofNW.center) to (NNofNW.center);
		\draw (NofNE.center) to (NNofNE.center);
\end{tikzpicture}}

\newcommand{\ladderab}{%
\begin{tikzpicture}[baseline=-1.5cm]
		\node[draw, rectangle] (coupon) at (-3.75, -1.5) {$\pi_{V_\mu}$};
		\node  (SW) at (-4, -2.5) {};
		\node  (SE) at (-3.5, -2.5) {};
		\node  (NW) at (-4, -1) {};
		\node  (NE) at (-3.5, -1) {};
		\node  (NofNW) at (-4, 0) {};
		\node  (NofNE) at (-3.5, 0) {};
		\node  (NNofNW) at (-4, .5) {};
		\node  (NNofNE) at (-3.5, .5) {};
		\draw (SW.center) to (coupon.230);
		\draw (SE.center) to (coupon.310);
		\draw (coupon.130) to (NW.center);
		\draw (coupon.50) to (NE.center);
		\draw (NW.center) to (NofNW.center);
		\draw (NE.center) to (NofNE.center);
		\draw [bend left=70, looseness=2] (NW.center) to (NofNE.center);
		\draw (NofNW.center) to (NNofNW.center);
		\draw (NofNE.center) to (NNofNE.center);
\end{tikzpicture}}

\newcommand{\ladderba}{%
\begin{tikzpicture}[baseline=-1.5cm]
		\node[draw, rectangle] (coupon) at (-3.75, -1.5) {$\pi_{V_\mu}$};
		\node  (SW) at (-4, -2.5) {};
		\node  (SE) at (-3.5, -2.5) {};
		\node  (NW) at (-4, -1) {};
		\node  (NE) at (-3.5, -1) {};
		\node  (NofNW) at (-4, 0) {};
		\node  (NofNE) at (-3.5, 0) {};
		\node  (NNofNW) at (-4, .5) {};
		\node  (NNofNE) at (-3.5, .5) {};
		\draw (SW.center) to (coupon.230);
		\draw (SE.center) to (coupon.310);
		\draw (coupon.130) to (NW.center);
		\draw (coupon.50) to (NE.center);
		\draw (NW.center) to (NofNW.center);
		\draw (NE.center) to (NofNE.center);
		\draw [bend left=70, looseness=2] (NE.center) to (NofNW.center);
		\draw (NofNW.center) to (NNofNW.center);
		\draw (NofNE.center) to (NNofNE.center);
\end{tikzpicture}}

\newcommand{\ladderbb}{%
\begin{tikzpicture}[baseline=-1.5cm]
		\node[draw, rectangle] (coupon) at (-3.75, -1.5) {$\pi_{V_\mu}$};
		\node  (SW) at (-4, -2.5) {};
		\node  (SE) at (-3.5, -2.5) {};
		\node  (NW) at (-4, -1) {};
		\node  (NE) at (-3.5, -1) {};
		\node  (NofNW) at (-4, 0) {};
		\node  (NofNE) at (-3.5, 0) {};
		\node  (NNofNW) at (-4, .5) {};
		\node  (NNofNE) at (-3.5, .5) {};
		\draw (SW.center) to (coupon.230);
		\draw (SE.center) to (coupon.310);
		\draw (coupon.130) to (NW.center);
		\draw (coupon.50) to (NE.center);
		\draw (NW.center) to (NofNW.center);
		\draw (NE.center) to (NofNE.center);
		\draw [bend left=90, looseness=2.5] (NE.center) to (NofNE.center);
		\draw (NofNW.center) to (NNofNW.center);
		\draw (NofNE.center) to (NNofNE.center);
\end{tikzpicture}}

\newcommand{\unitcap}{%
\begin{tikzpicture}[baseline=.5cm]
	\node (W) at (-.5, 0) {$\bullet$};
	\node (E) at (.5, 0) {$\bullet$};
	\draw [bend left=90, looseness=2] (W.center) to (E.center);
\end{tikzpicture}}

\newcommand{\twocap}{%
\begin{tikzpicture}[baseline= .5cm]
	\node (1) at (-1, 0) {};
	\node (2) at (-.25, 0) {};
	\node (3) at (.25, 0) {};
	\node (4) at (1, 0) {};
	\draw [bend left=90, looseness=2] (1) to (2);
	\draw [bend left=90, looseness=2] (3) to (4);
\end{tikzpicture}}

\newcommand{\HurwitzTerma}{%
\begin{tikzpicture}[baseline=.5cm]
	\node (a) at (-1.5, 0) {};
	\node (b) at (-.5, 0) {};
	\node (c) at (.5, 0) {};
	\node (d) at (1.5, 0) {};
	\node [draw, rectangle] (1) at (-.5,1) {$\mu$};
	\node [draw, rectangle] (2) at (.5,1) {$\mu$};
	\draw (a) to (1.250);
	\draw (c) to (1.290);
	\draw (b) to (2.250);
	\draw (d) to (2.290);
	\draw [bend left=90, looseness=2] (1) to (2);
\end{tikzpicture}}

\newcommand{\HurwitzTermb}{%
\begin{tikzpicture}[baseline=.5cm]
	\node (a) at (-1.5, 0) {};
	\node (b) at (-.5, 0) {};
	\node (c) at (.5, 0) {};
	\node (d) at (1.5, 0) {};
	\node [draw, rectangle] (1) at (-.5,1) {$\mu$};
	\node [draw, rectangle] (2) at (.5,1) {$\mu$};
	\draw (a) to (1.250);
	\draw (d) to (1.290);
	\draw (b) to (2.250);
	\draw (c) to (2.290);
	\draw [bend left=90, looseness=2] (1) to (2);
\end{tikzpicture}}


\newcommand{\TracialJordanAxiom}{%
\begin{tikzpicture}[baseline=2cm]
	\node [text height=1.5ex] (a) at (-.5,0) {$a$};
	\node [text height=1.5ex] (b) at (0,0) {$b$};
	\node [text height=1.5ex] (c) at (.5,0) {$c$};
	\node [draw, rectangle] (sym) at  (0,1) {$\; \Sym^3 \;$};
	\node [draw, rectangle] (wedge) at (-.25,3) {$\; \wedge^2 \;$};
	\node [text height=1.5ex] (y) at (-.5,4) {$y$};
	\node (output) at (0,4) {};
	\node (r) at (-.5, 2) {};
	\node (s) at (0,2) {};
	\node (t) at (.25,1.75) {};
	\draw (a) to (sym.216);
	\draw (b) to (sym.270);
	\draw (c) to (sym.324);
	\draw (r.center) to (sym.144);
	\draw (r.center) to (s.center);
	\draw (s.center) to (t.center);
	\draw (sym.90) to (t.center);
	\draw (sym.36) to (t.center);
	\draw (r.center) to (wedge.233);
	\draw (s.center) to (wedge.307);
	\draw (wedge.127) to (y);
	\draw (wedge.53) to (output);
\end{tikzpicture}}

\newcommand{\TracialJordanAxiomLonga}{%
\begin{tikzpicture}[baseline=2cm]
	\node [text height=1.5ex] (a) at (-.5,0) {$a$};
	\node [text height=1.5ex] (b) at (0,0) {$b$};
	\node [text height=1.5ex] (c) at (.5,0) {$c$};
	\node [draw, rectangle] (sym) at  (0,1) {$\; \Sym^3 \;$};
	\node [text height=1.5ex] (y) at (-.5,4) {$y$};
	\node (output) at (0,4) {};
	\node (r) at (-.5, 2) {};
	\node (s) at (0,2) {};
	\node (t) at (.25,1.75) {};
	\draw (a) to (sym.216);
	\draw (b) to (sym.270);
	\draw (c) to (sym.324);
	\draw (r.center) to (sym.144);
	\draw (r.center) to (s.center);
	\draw (s.center) to (t.center);
	\draw (sym.90) to (t.center);
	\draw (sym.36) to (t.center);
	\draw (r.center) to (y);
	\draw (s.center) to (output);
\end{tikzpicture}}

\newcommand{\TracialJordanAxiomLongb}{%
\begin{tikzpicture}[baseline=2cm]
	\node [text height=1.5ex] (a) at (-.5,0) {$a$};
	\node [text height=1.5ex] (b) at (0,0) {$b$};
	\node [text height=1.5ex] (c) at (.5,0) {$c$};
	\node [draw, rectangle] (sym) at  (0,1) {$\; \Sym^3 \;$};
	\node [text height=1.5ex] (y) at (-.5,4) {$y$};
	\node (output) at (0,4) {};
	\node (r) at (-.5, 2) {};
	\node (s) at (0,2) {};
	\node (t) at (.25,1.75) {};
	\draw (a) to (sym.216);
	\draw (b) to (sym.270);
	\draw (c) to (sym.324);
	\draw (r.center) to (sym.144);
	\draw (r.center) to (s.center);
	\draw (s.center) to (t.center);
	\draw (sym.90) to (t.center);
	\draw (sym.36) to (t.center);
	\draw (r.center) to (output);
	\draw (s.center) to (y);
\end{tikzpicture}}

	
\newcommand{\TracelessJordanAxioma}{%
\begin{tikzpicture}[baseline=2cm]
	\node (a) at (-.5,0) {};
	\node (b) at (0,0) {};
	\node (c) at (.5,0) {};
	\node [draw, rectangle] (sym) at  (0,1) {$\; \Sym^3 \;$};
	\node [draw, rectangle] (wedge) at (-.25,3) {$\; \wedge^2 \;$};
	\node (y) at (-.5,4) {};
	\node (output) at (0,4) {};
	\node (r) at (-.5, 2) {};
	\node (s) at (0,2) {};
	\node (t) at (.25,1.75) {};
	\draw (a) to (sym.216);
	\draw (b) to (sym.270);
	\draw (c) to (sym.324);
	\draw (r.center) to (sym.144);
	\draw (r.center) to (s.center);
	\draw (s.center) to (t.center);
	\draw (sym.90) to (t.center);
	\draw (sym.36) to (t.center);
	\draw (r.center) to (wedge.233);
	\draw (s.center) to (wedge.307);
	\draw (wedge.127) to (y);
	\draw (wedge.53) to (output);
\end{tikzpicture}}

\newcommand{\TracelessJordanAxiomb}{%
\begin{tikzpicture}[baseline=2cm]
	\node (a) at (-.5,0) {};
	\node (b) at (0,0) {};
	\node (c) at (.5,0) {};
	\node [draw, rectangle] (sym) at  (0,1) {$\; \Sym^3 \;$};
	\node [draw, rectangle] (wedge) at (-.25,3) {$\; \wedge^2 \;$};
	\node (y) at (-.5,4) {};
	\node (output) at (0,4) {};
	\node (s) at (0,2) {};
	\draw (a) to (sym.216);
	\draw (b) to (sym.270);
	\draw (c) to (sym.324);
	\draw (wedge.233) to (sym.144);
	\draw (sym.90) to (s.center);
	\draw (sym.36) to (s.center);
	\draw (s.center) to (wedge.307);
	\draw (wedge.125) to (y);
	\draw (wedge.53) to (output);
\end{tikzpicture}}

\newcommand{\TracialThreefoldCHa}{%
\begin{tikzpicture}[baseline=1.5cm]
	\node (a) at (-.5,0) {};
	\node (b) at (0,0) {};
	\node (c) at (.5,0) {};
	\node [draw, rectangle] (sym) at  (0,1) {$\; \Sym^3 \;$};
	\node (d) at (0,3) {};
	\node (s) at (-.25,1.75) {};
	\node (t) at (0,2.25) {};
	\draw (a) to (sym.216);
	\draw (b) to (sym.270);
	\draw (c) to (sym.324);
	\draw (sym.144) to (s.center);
	\draw (sym.90) to (s.center);
	\draw (s.center) to (t.center);
	\draw (sym.35) to (t.center);
	\draw (t.center) to (d);
\end{tikzpicture}}

\newcommand{\TracialThreefoldCHb}{%
\begin{tikzpicture}[baseline=1.5cm]
	\node (a) at (-.5,0) {};
	\node (b) at (0,0) {};
	\node (c) at (.5,0) {};
	\node [draw, rectangle] (sym) at  (0,1) {$\; \Sym^3 \;$};
	\node (d) at (0,3) {};
	\node (s) at (-.25,2) {};
	\node (u) at (.5,2) {$\bullet$};
	\draw (a) to (sym.216);
	\draw (b) to (sym.270);
	\draw (c) to (sym.324);
	\draw (sym.144) to (s.center);
	\draw (sym.90) to (s.center);
	\draw (s.center) to (d);
	\draw (sym.36) to (u.center);
\end{tikzpicture}}

\newcommand{\TracialThreefoldCHc}{%
\begin{tikzpicture}[baseline=1.5cm]
	\node (a) at (-.5,0) {};
	\node (b) at (0,0) {};
	\node (c) at (.5,0) {};
	\node [draw, rectangle] (sym) at  (0,1) {$\; \Sym^3 \;$};
	\node (d) at (0,3) {};
	\node (u2) at (0,2) {$\bullet$};
	\node (u3) at (.5,2) {$\bullet$};
	\draw (a) to (sym.216);
	\draw (b) to (sym.270);
	\draw (c) to (sym.324);
	\draw (sym.144) to (d);
	\draw (sym.90) to (u2.center);
	\draw (sym.36) to (u3.center);	
\end{tikzpicture}}

\newcommand{\TracialThreefoldCHd}{%
\begin{tikzpicture}[baseline=1.5cm]
	\node (a) at (-.5,0) {};
	\node (b) at (0,0) {};
	\node (c) at (.5,0) {};
	\node [draw, rectangle] (sym) at  (0,1) {$\; \Sym^3 \;$};
	\node (d) at (0,3) {};
	\draw (a) to (sym.216);
	\draw (b) to (sym.270);
	\draw (c) to (sym.324);
	\draw (sym.144) to (d);
	\draw [bend left=90, looseness=2] (sym.90) to (sym.36);
\end{tikzpicture}}

\newcommand{\TracialThreefoldCHe}{%
\begin{tikzpicture}[baseline=1.5cm]
	\node (a) at (-.5,0) {};
	\node (b) at (0,0) {};
	\node (c) at (.5,0) {};
	\node [draw, rectangle] (sym) at  (0,1) {$\; \Sym^3 \;$};
	\node (d) at (0,3) {};
	\node (u) at (0,2.5) {$\bullet$};
	\node (s) at (0,2) {};
	\draw (a) to (sym.216);
	\draw (b) to (sym.270);
	\draw (c) to (sym.324);
	\draw (sym.144) to (s.center);
	\draw (sym.90) to (s.center);
	\draw (sym.36) to (s.center);
	\draw (u.center) to (d);	
\end{tikzpicture}}

\newcommand{\TracialThreefoldCHf}{%
\begin{tikzpicture}[baseline=1.5cm]
	\node (a) at (-.5,0) {};
	\node (b) at (0,0) {};
	\node (c) at (.5,0) {};
	\node [draw, rectangle] (sym) at  (0,1) {$\; \Sym^3 \;$};
	\node (d) at (0,3) {};
	\node (u1) at (-.5,1.75) {$\bullet$};
	\node (u2) at (0,2.25) {$\bullet$};	
	\draw (a) to (sym.216);
	\draw (b) to (sym.270);
	\draw (c) to (sym.324);
	\draw (sym.144) to (u1.center);
	\draw (u2.center) to (d);	
	\draw [bend left=90, looseness=2] (sym.90) to (sym.36);
\end{tikzpicture}}

\newcommand{\TracialThreefoldCHg}{%
\begin{tikzpicture}[baseline=1.5cm]
	\node (a) at (-.5,0) {};
	\node (b) at (0,0) {};
	\node (c) at (.5,0) {};
	\node [draw, rectangle] (sym) at  (0,1) {$\; \Sym^3 \;$};
	\node (d) at (0,3) {};
	\node (u1) at (-.5,1.75) {$\bullet$};
	\node (u2) at (0,1.75) {$\bullet$};
	\node (u3) at (.5,1.75) {$\bullet$};
	\node (u4) at (0,2.25) {$\bullet$};
	\draw (a) to (sym.216);
	\draw (b) to (sym.270);
	\draw (c) to (sym.324);
	\draw (sym.144) to (u1.center);
	\draw (sym.90) to (u2.center);
	\draw (sym.36) to (u3.center);
	\draw (u4.center) to (d);	
\end{tikzpicture}}


\newcommand{\firstthreebox}{%
\begin{tikzpicture}[baseline=-0.5ex,scale=0.8]
	\node (a) at (150:.8cm) {};
	\node (b) at (30:.8cm) {};
	\node (c) at (270:.8cm) {};
	\node (a') at (150:.2cm) {$\bullet$};
	\node (b') at (30:.2cm) {$\bullet$};
	\node (c') at (270:.2cm) {$\bullet$};
	\draw (a.center) to (a'.center);
	\draw (b.center) to (b'.center);
	\draw (c.center) to (c'.center);
\end{tikzpicture}}

\newcommand{\secondthreebox}{%
\begin{tikzpicture}[baseline=-0.5ex,scale=0.8]
	\node (a) at (150:.8cm) {};
	\node (b) at (30:.8cm) {};
	\node (c) at (270:.8cm) {};
	\node (a') at (150:.2cm) {$\bullet$};
	\draw (a.center) to (a'.center);
	\draw [bend right=45, looseness=.5] (b.center) to (c.center);
\end{tikzpicture}}

\newcommand{\thirdthreebox}{%
\begin{tikzpicture}[baseline=-0.5ex,scale=0.8]
	\node (a) at (150:.8cm) {};
	\node (b) at (30:.8cm) {};
	\node (c) at (270:.8cm) {};
	\node (b') at (30:.2cm) {$\bullet$};
	\draw [bend left=45, looseness=.5] (a.center) to (c.center);
	\draw (b.center) to (b'.center);
\end{tikzpicture}}

\newcommand{\fourththreebox}{%
\begin{tikzpicture}[baseline=-0.5ex,scale=0.8]
	\node (a) at (150:.8cm) {};
	\node (b) at (30:.8cm) {};
	\node (c) at (270:.8cm) {};
	\node (c') at (270:.2cm) {$\bullet$};
	\draw [bend right=45, looseness=1] (a.center) to (b.center);
	\draw (c.center) to (c'.center);
\end{tikzpicture}}


\newcommand{\dogbone}{%
\begin{tikzpicture}[baseline=0cm]
	\node (a) at (-.5,0) {$\bullet$};
	\node (b) at (.5,0) {$\bullet$};
	\draw (a.center) to (b.center);
\end{tikzpicture}}

\newcommand{\counitunit}{%
\begin{tikzpicture}[baseline=-0.5ex]
    \node (a) at (0,.15) {$\bullet$};
    \node (b) at (0,-.15) {$\bullet$};
    \draw (a.center) -- ++(0,.45);
    \draw (b.center) -- ++(0,-.45);
\end{tikzpicture}%
}

\newcommand{\onevalent}{%
\begin{tikzpicture}[baseline=0cm]
	\node (a) at (0,-.5) {};
	\node (b) at (0,.5) {$\bullet$};
	\draw (a.center) to (b.center);
\end{tikzpicture}}

\newcommand{\leftunitvertex}{%
\begin{tikzpicture}[baseline=0cm]
	\node (a) at (-.5,.5) {$\bullet$};
	\node (b) at (.5,.5) {};
	\node (c) at (0,-.5) {};
	\node (d) at (0,0) {};
	\draw (a.center) to (d.center);
	\draw (b.center) to (d.center);
	\draw (c.center) to (d.center);
\end{tikzpicture}}

\newcommand{\rightunitvertex}{%
\begin{tikzpicture}[baseline=0cm]
	\node (a) at (-.5,.5) {};
	\node (b) at (.5,.5) {$\bullet$};
	\node (c) at (0,-.5) {};
	\node (d) at (0,0) {};
	\draw (a.center) to (d.center);
	\draw (b.center) to (d.center);
	\draw (c.center) to (d.center);
\end{tikzpicture}}

\newcommand{\TracelessJordanProofa}{%
\begin{tikzpicture}[baseline=2cm]
	\node (a) at (-.5,0) {};
	\node (b) at (0,0) {};
	\node (c) at (.5,0) {};
	\node [draw, rectangle] (sym) at  (0,1) {$\; \Sym^3 \;$};
	\node [draw, rectangle] (wedge) at (-.25,3) {$\; \wedge^2 \;$};
	\node (y) at (-.5,4) {};
	\node (output) at (0,4) {};
	\node (s) at (-.5,2) {};
	\node (t) at (0,1.75) {};
	\node (u) at (.5,2) {};
	\draw (a) to (sym.216);
	\draw (b) to (sym.270);
	\draw (c) to (sym.324);
	\draw (sym.144) to (s.center);
	\draw (sym.90) to (t.center);
	\draw (sym.36) to (u.center);
	\draw (s.center) to (wedge.233);
	\draw (u.center) to (wedge.307);
	\draw (wedge.125) to (y);
	\draw (wedge.53) to (output);
	\draw (s.center) to (t.center);
	\draw (t.center) to (u.center);
\end{tikzpicture}}

\newcommand{\TracelessJordanProofb}{%
\begin{tikzpicture}[baseline=2cm]
	\node (a) at (-.5,0) {};
	\node (b) at (0,0) {};
	\node (c) at (.5,0) {};
	\node [draw, rectangle] (sym) at  (0,1) {$\; \Sym^3 \;$};
	\node [draw, rectangle] (wedge) at (-.25,3) {$\; \wedge^2 \;$};
	\node (y) at (-.5,4) {};
	\node (output) at (0,4) {};
	\node (s) at (-.5,2) {};
	\draw (a) to (sym.216);
	\draw (b) to (sym.270);
	\draw (c) to (sym.324);
	\draw (wedge.233) to (sym.144);
	\draw [bend left=90, looseness=2] (sym.90) to (sym.36);
	\draw (s.center) to (wedge.307);
	\draw (wedge.125) to (y);
	\draw (wedge.53) to (output);
\end{tikzpicture}}

\newcommand{\TracelessJordanProofc}{%
\begin{tikzpicture}[baseline=2cm]
	\node (a) at (-.5,0) {};
	\node (b) at (0,0) {};
	\node (c) at (.5,0) {};
	\node [draw, rectangle] (sym) at  (0,1) {$\; \Sym^3 \;$};
	\node [draw, rectangle] (wedge) at (-.25,3) {$\; \wedge^2 \;$};
	\node (y) at (-.5,4) {};
	\node (output) at (0,4) {};
	\node (s) at (-.5,2) {};
	\draw (a) to (sym.216);
	\draw (b) to (sym.270);
	\draw (c) to (sym.324);
	\draw (wedge.233) to (sym.144);
	\draw (sym.90) to (s.center);
	\draw (sym.36) to (wedge.307);
	\draw (wedge.125) to (y);
	\draw (wedge.53) to (output);
\end{tikzpicture}}

\newcommand{\capadjointb}{%
\begin{tikzpicture}[baseline=.5cm]
	\node (a) at (-1.5,0) {};
	\node (b) at (-.5,0) {};
	\node (c) at (.5,1) {};
	\node (d) at (1.5,1) {};
	\draw (c.center) to [out=-90, in = -90, looseness =2] (d.center);
	\draw (a.center) to [out = 90, in = 90, looseness = 2] (c.center);
	\draw (b.center) to [out = 90, in = 90, looseness = 2] (d.center);
	\pgfresetboundingbox
	\draw(-1.5,0) (1.5,2);
\end{tikzpicture}}


\newcommand{\trivalentadjointa}{%
\begin{tikzpicture}[baseline=.5cm]
	\node (a) at (-1,0) {};
	\node (b) at (0,0) {};
	\node (c) at (1,0) {};
	\node (t) at (-.5,.5) {};
	\node (c') at (1,.5) {};
	\draw (a.center) to (t.center);
	\draw (b.center) to (t.center);
	\draw (c.center) to (c'.center);
	\draw (t.center) to [bend left=90, looseness=2] (c'.center);
	\pgfresetboundingbox
	\draw(-1,0) (1.25,1.5);
\end{tikzpicture}}

\newcommand{\trivalentadjointb}{%
\begin{tikzpicture}[baseline=.5cm]
	\node (a) at (-1,0) {};
	\node (b) at (0,0) {};
	\node (c) at (1,0) {};
	\node (t) at (1,.5) {};
	\draw (c.center) to (t.center);
	\draw (a.center) to [out=90,in=135, looseness =1] (t.center);
	\draw (b.center) to [out=90, in=45, looseness = 3] (t.center);
	\pgfresetboundingbox
	\draw(-1,0) (1.25,1.25);
\end{tikzpicture}}

\newcommand{\trivalentadjointc}{%
\begin{tikzpicture}[baseline=.5cm]
	\node (a) at (-1,0) {};
	\node (b) at (0,0) {};
	\node (c) at (1,0) {};
	\node (t) at (1,.5) {};
        \draw (c.center) to (t.center);
        \draw (a.center) to [out=90,in=45,looseness=2] (t.center);
        \draw (b.center) to [out=90,in=135,looseness=1] (t.center);
	\pgfresetboundingbox
	\draw(-1,0) (1.25,1.25);      
\end{tikzpicture}}

\newcommand{\JordanLowerTermsa}{%
\begin{tikzpicture}[baseline=2cm]
	\node (a) at (-.5,0) {};
	\node (b) at (.5,0) {};
	\node [draw, rectangle, minimum width = 1.4cm, minimum height = .7cm] (sym) at  (0,1) {$\; \Sym^2 \;$};
	\node [draw, rectangle, minimum width = 1.4cm, minimum height = .7cm] (wedge) at (0,3) {$\; \wedge^2 \;$};
	\node (y) at (-.5,4) {};
	\node (output) at (.5,4) {};
	\node (s) at (0,1.75) {};
	\node (t) at (0,2.25) {};
	\draw (a) to (sym.216);
	\draw (b) to (sym.324);
	\draw (sym.144) to (s.center);
	\draw (sym.36) to (s.center);
	\draw (t.center) to (wedge.216);
	\draw (t.center) to (wedge.324);
	\draw (wedge.144) to (y);
	\draw (wedge.36) to (output);
	\draw (s.center) to (t.center);
\end{tikzpicture}}

\newcommand{\JordanLowerTermsb}{%
\begin{tikzpicture}[baseline=2cm]
	\node (a) at (-.5,0) {};
	\node (b) at (.5,0) {};
	\node [draw, rectangle, minimum width = 1.4cm, minimum height = .7cm] (sym) at  (0,1) {$\; \Sym^2 \;$};
	\node [draw, rectangle, minimum width = 1.4cm, minimum height = .7cm] (wedge) at (0,3) {$\; \wedge^2 \;$};
	\node (y) at (-.5,4) {};
	\node (output) at (.5,4) {};
	\node (s) at (-.5,2) {};
	\node (t) at (+.5,2) {};
	\draw (a) to (sym.216);
	\draw (b) to (sym.324);
	\draw (sym.144) to (s.center);
	\draw (sym.36) to (t.center);
	\draw (s.center) to (wedge.216);
	\draw (t.center) to (wedge.324);
	\draw (wedge.144) to (y);
	\draw (wedge.36) to (output);
	\draw (s.center) to (t.center);
\end{tikzpicture}}

\newcommand{\JordanLowerTermsc}{%
\begin{tikzpicture}[baseline=2cm]
	\node (a) at (-.5,0) {};
	\node (b) at (.5,0) {};
	\node [draw, rectangle, minimum width = 1.4cm, minimum height = .7cm] (sym) at  (0,1) {$\; \Sym^2 \;$};
	\node [draw, rectangle, minimum width = 1.4cm, minimum height = .7cm] (wedge) at (0,3) {$\; \wedge^2 \;$};
	\node (y) at (-.5,4) {};
	\node (output) at (.5,4) {};
	\draw (a) to (sym.216);
	\draw (b) to (sym.324);
	\draw (sym.144) to (wedge.216);
	\draw (sym.36) to (wedge.324);
	\draw (wedge.144) to (y);
	\draw (wedge.36) to (output);
\end{tikzpicture}}

% \newfig{\pentagon}{\pentagonbox}{
% \begin{tikzpicture}[use Hobby shortcut]
%         \coordinate (E1) at (90
% \end{tikzpicture}
% }

%%%%%%%%%%%

\tikzset{
  Pcoupon/.style={
    rectangle,
    draw,
    fill=white,
    inner sep=0pt,
    outer sep=0pt,
    minimum width=4.5mm,
    minimum height=3mm,
    font=\scriptsize
  },
  unitvertex/.style={
    circle,
    fill,
    inner sep=0pt,
    minimum size=1.4mm
  }
}

% #1 and #2 describe the left and right internal edges:
% 0 = identity, 1 = P, 2 = unit-counit.
\newcommand{\PBigon}[2]{%
\begin{tikzpicture}[
    baseline={(0,.60)},
    line width=.45pt,
    x=.9cm,
    y=.9cm
  ]
  % Lower boundary P
  \draw (0,-.55) -- (0,-.47);
  \draw (0,-.13) -- (0,0);
  \node[Pcoupon] at (0,-.30) {$P$};

  % Upper boundary P
  \draw (0,1.20) -- (0,1.33);
  \draw (0,1.67) -- (0,1.75);
  \node[Pcoupon] at (0,1.50) {$P$};

  % Left internal edge
  \ifcase#1
    \draw (0,0)
      .. controls (-.36,.23) and (-.36,.97) ..
      (0,1.20);
  \or
    \draw (0,0)
      .. controls (-.29,.17) and (-.35,.34) ..
      (-.34,.43);
    \draw (-.34,.77)
      .. controls (-.35,.86) and (-.29,1.03) ..
      (0,1.20);
    \node[Pcoupon] at (-.34,.60) {$P$};
  \or
    \draw (0,0)
      .. controls (-.29,.17) and (-.34,.36) ..
      (-.33,.48);
    \node[unitvertex] at (-.33,.48) {};
    \node[unitvertex] at (-.33,.72) {};
    \draw (-.33,.72)
      .. controls (-.34,.84) and (-.29,1.03) ..
      (0,1.20);
  \fi

  % Right internal edge
  \ifcase#2
    \draw (0,0)
      .. controls (.36,.23) and (.36,.97) ..
      (0,1.20);
  \or
    \draw (0,0)
      .. controls (.29,.17) and (.35,.34) ..
      (.34,.43);
    \draw (.34,.77)
      .. controls (.35,.86) and (.29,1.03) ..
      (0,1.20);
    \node[Pcoupon] at (.34,.60) {$P$};
  \or
    \draw (0,0)
      .. controls (.29,.17) and (.34,.36) ..
      (.33,.48);
    \node[unitvertex] at (.33,.48) {};
    \node[unitvertex] at (.33,.72) {};
    \draw (.33,.72)
      .. controls (.34,.84) and (.29,1.03) ..
      (0,1.20);
  \fi
\end{tikzpicture}%
}

\newcommand{\PIdentity}{%
\begin{tikzpicture}[
    baseline={(0,.60)},
    line width=.45pt,
    x=.9cm,
    y=.9cm
  ]
  \draw (0,-.10) -- (0,.43);
  \draw (0,.77) -- (0,1.30);
  \node[Pcoupon] at (0,.60) {$P$};
\end{tikzpicture}%
}

\newcommand{\twistone}{%
\begin{tikzpicture}[
    baseline=(current bounding box.center),
    scale=1.5,
    every path/.style={
        draw,
        line width=.45pt,
        line cap=round,
        line join=round
    }
]
    % Under-strand, curving smoothly into the right vertex
    \draw
        (0,0)
        -- (.34,.34)
        .. controls (.43,.43) and (.46,.56) ..
        (.375,.625);

    % Over-strand, curving smoothly into the left vertex
    \draw[
        preaction={draw=white,line width=2.2pt}
    ]
        (.5,0)
        -- (.16,.34)
        .. controls (.07,.43) and (.04,.56) ..
        (.125,.625);

    % Upper outgoing strands
    \draw
        (.125,.625)
        .. controls (.07,.66) and (.025,.71) ..
        (0,.75);

    \draw
        (.375,.625)
        .. controls (.43,.66) and (.475,.71) ..
        (.5,.75);

    % Internal edge
    \draw
        (.125,.625)
        -- (.375,.625);
\end{tikzpicture}%
}

\newcommand{\twisttwo}{%
\begin{tikzpicture}[
    baseline=(current bounding box.center),
    scale=1.5,
    every path/.style={
        draw,
        line width=.45pt,
        line cap=round,
        line join=round
    }
]
    % Under-strand, curving smoothly into the lower vertex
    \draw
        (.25,.25)
        -- (-.14,-.14)
        .. controls (-.23,-.23) and (-.32,-.21) ..
        (-.375,-.125);

    % Over-strand, curving smoothly into the upper vertex
    \draw[
        preaction={draw=white,line width=2.2pt}
    ]
        (.25,-.25)
        -- (-.14,.14)
        .. controls (-.23,.23) and (-.32,.21) ..
        (-.375,.125);

    % Left outgoing strands
    \draw
        (-.375,.125)
        .. controls (-.42,.18) and (-.47,.23) ..
        (-.5,.25);

    \draw
        (-.375,-.125)
        .. controls (-.42,-.18) and (-.47,-.23) ..
        (-.5,-.25);

    % Internal edge
    \draw
        (-.375,-.125)
        -- (-.375,.125);
\end{tikzpicture}%
}